% ─────────────────────────────────────────────────────────────────────────────
%  presentation.tex  –  Presentación Beamer
%  Compilar: pdflatex presentation.tex  (×2)
% ─────────────────────────────────────────────────────────────────────────────
\documentclass[aspectratio=169,12pt]{beamer}

\usetheme{Madrid}
\usecolortheme{whale}
\usefonttheme{professionalfonts}

% ── Colores ──────────────────────────────────────────────────────────────────
\definecolor{primary}{RGB}{0,62,126}
\definecolor{accent}{RGB}{0,153,76}
\definecolor{warn}{RGB}{200,30,30}
\definecolor{codebg}{RGB}{248,248,248}
\definecolor{codegreen}{RGB}{34,139,34}
\definecolor{codegray}{RGB}{128,128,128}
\definecolor{codepurple}{RGB}{148,0,211}
\definecolor{codeblue}{RGB}{0,51,178}

\setbeamercolor{frametitle}{bg=primary,fg=white}
\setbeamercolor{title}{bg=primary,fg=white}
\setbeamercolor{structure}{fg=primary}
\setbeamercolor{block title}{bg=primary!80,fg=white}
\setbeamercolor{block body}{bg=primary!8}
\setbeamercolor{block title alerted}{bg=warn!80,fg=white}
\setbeamercolor{block body alerted}{bg=warn!8}
\setbeamercolor{block title example}{bg=accent!80,fg=white}
\setbeamercolor{block body example}{bg=accent!8}

% ── Paquetes ─────────────────────────────────────────────────────────────────
\usepackage[utf8]{inputenc}
\usepackage[T1]{fontenc}
\usepackage[spanish,es-noshorthands]{babel}
\usepackage{lmodern}
\usepackage{amsmath,amssymb,amsthm}
\usepackage{listings}
\usepackage{xcolor}
\usepackage{tikz}
\usetikzlibrary{shapes.geometric,arrows.meta,positioning,fit,calc}
\usepackage{array}
\usepackage{multicol}

% ── Listings ─────────────────────────────────────────────────────────────────
\lstdefinestyle{javaslide}{
  backgroundcolor=\color{codebg},
  commentstyle=\color{codegreen}\itshape,
  keywordstyle=\color{codeblue}\bfseries,
  stringstyle=\color{codepurple},
  basicstyle=\ttfamily\tiny,
  breaklines=true,
  keepspaces=true,
  numbers=left,
  numberstyle=\tiny\color{codegray},
  numbersep=3pt,
  tabsize=2,
  language=Java,
  morekeywords={result},
  frame=single,
  rulecolor=\color{gray!40},
}
\lstdefinestyle{cppslide}{
  backgroundcolor=\color{codebg},
  commentstyle=\color{codegreen}\itshape,
  keywordstyle=\color{codeblue}\bfseries,
  basicstyle=\ttfamily\tiny,
  breaklines=true,
  numbers=none,
  tabsize=2,
  frame=single,
  rulecolor=\color{gray!40},
  language=C++,
}
\lstdefinestyle{asmslide}{
  backgroundcolor=\color{codebg},
  commentstyle=\color{codegreen}\itshape,
  basicstyle=\ttfamily\tiny,
  breaklines=true,
  numbers=none,
  frame=single,
  rulecolor=\color{gray!40},
}
\lstdefinestyle{smtslide}{
  backgroundcolor=\color{codebg},
  commentstyle=\color{codegreen}\itshape,
  keywordstyle=\color{codepurple}\bfseries,
  basicstyle=\ttfamily\tiny,
  breaklines=true,
  numbers=left,
  numberstyle=\tiny\color{codegray},
  numbersep=3pt,
  frame=single,
  rulecolor=\color{gray!40},
  morekeywords={declare-fun,assert,check-sat,push,pop,not,and,or,Int},
}

% ── Macros ───────────────────────────────────────────────────────────────────
\newcommand{\WP}[2]{\mathrm{WP}(#1,\,#2)}
\newcommand{\hoare}[3]{\{#1\}\;#2\;\{#3\}}
\newcommand{\subst}[3]{#1[#3/#2]}
\newcommand{\reg}[1]{\texttt{\%#1}}

% ── Footer ───────────────────────────────────────────────────────────────────
\setbeamertemplate{footline}{%
  \leavevmode
  \hbox{%
    \begin{beamercolorbox}[wd=.333\paperwidth,ht=2.25ex,dp=1ex,center]{author in head/foot}%
      \usebeamerfont{author in head/foot}\insertshortauthor
    \end{beamercolorbox}%
    \begin{beamercolorbox}[wd=.334\paperwidth,ht=2.25ex,dp=1ex,center]{title in head/foot}%
      \usebeamerfont{title in head/foot}\insertshorttitle
    \end{beamercolorbox}%
    \begin{beamercolorbox}[wd=.333\paperwidth,ht=2.25ex,dp=1ex,right]{date in head/foot}%
      \usebeamerfont{date in head/foot}
      \insertframenumber{} / \inserttotalframenumber\hspace*{1em}
    \end{beamercolorbox}}%
  \vskip0pt%
}
\setbeamertemplate{navigation symbols}{}

% ── Metadatos ────────────────────────────────────────────────────────────────
\title[Compilador JavaSubset $\to$ x86-64]{%
  Compilador \texttt{JavaSubset} $\rightarrow$ x86-64\\[0.3em]
  {\large con Verificaci\'on Formal (Hoare + Z3)}}
\subtitle{}
\author{Jaime Hernández}
\institute{Dise\~no de Compiladores}
\date{\today}

% ═════════════════════════════════════════════════════════════════════════════
\begin{document}

% ─── Portada ─────────────────────────────────────────────────────────────────
\begin{frame}[plain]
  \titlepage
\end{frame}

% ─── Índice ──────────────────────────────────────────────────────────────────
\begin{frame}{\'Indice}
  \begin{multicols}{2}
    \tableofcontents
  \end{multicols}
\end{frame}

% ═════════════════════════════════════════════════════════════════════════════
\section{Arquitectura general}
% ─────────────────────────────────────────────────────────────────────────────

\begin{frame}{Pipeline del compilador}
\centering
\begin{tikzpicture}[
  scale=0.78, transform shape,
  box/.style={rectangle, draw, rounded corners=2pt,
              minimum width=2.4cm, minimum height=0.55cm,
              align=center, font=\scriptsize\ttfamily, fill=blue!10},
  ver/.style={rectangle, draw, rounded corners=2pt,
              minimum width=2.4cm, minimum height=0.55cm,
              align=center, font=\scriptsize\ttfamily, fill=accent!15},
  arr/.style={-Stealth, thick},
  lbl/.style={font=\tiny\itshape, text=gray!80}
]
\node[box] (src) {Fuente .java};
\node[box, below=0.35cm of src]  (lex) {Lexer (Flex)};
\node[box, below=0.35cm of lex]  (par) {Parser (Bison)};
\node[box, below=0.35cm of par]  (sem) {Sem\'antico};
\node[box, below=0.35cm of sem]  (ir)  {Gen. IR};
\node[box, below=0.35cm of ir]   (opt) {Optim. IR};
\node[box, below=0.35cm of opt]  (reg) {Reg. Alloc};
\node[box, below=0.35cm of reg]  (x86) {x86-64 Gen};
\node[box, below=0.35cm of x86]  (asm) {Ensamblador .s};

\draw[arr] (src)--(lex);
\draw[arr] (lex)--(par) node[lbl,right,midway]{tokens};
\draw[arr] (par)--(sem) node[lbl,right,midway]{AST};
\draw[arr] (sem)--(ir)  node[lbl,right,midway]{tipos};
\draw[arr] (ir)--(opt)  node[lbl,right,midway]{IR};
\draw[arr] (opt)--(reg) node[lbl,right,midway]{IR opt.};
\draw[arr] (reg)--(x86) node[lbl,right,midway]{RegMap};
\draw[arr] (x86)--(asm) node[lbl,right,midway]{AT\&T};

\node[ver, right=2.8cm of opt, align=center] (wpc) {VCGenerator\\WPCalculus};
\node[ver, below=0.35cm of wpc, align=center] (smt) {SMTExporter\\(.smt2)};
\node[ver, below=0.35cm of smt] (z3)  {Z3 Solver};

\draw[arr, dashed, accent!80]
  (sem) -| node[lbl,above,near end]{--verify} (wpc);
\draw[arr, accent!80] (wpc)--(smt);
\draw[arr, accent!80] (smt)--(z3);
\end{tikzpicture}
\end{frame}

% ─────────────────────────────────────────────────────────────────────────────
\begin{frame}{Subconjunto Java soportado}
\begin{columns}[t]
  \column{0.48\textwidth}
  \begin{block}{Construcci\'ones}
    \begin{itemize}\small
      \item Tipos: \texttt{int}, \texttt{double}, \texttt{boolean}
      \item Variables locales + m\'etodos est\'aticos
      \item Asignaci\'on, aritm\'etica, comparaci\'on
      \item \texttt{if}/\texttt{else}, \texttt{while}
      \item \texttt{return}, llamadas a m\'etodos
      \item \texttt{System.out.println}
    \end{itemize}
  \end{block}

  \column{0.52\textwidth}
  \begin{alertblock}{Anotaciones de verificaci\'on (Fase 5)}
    \begin{tabular}{ll}\small
      \texttt{@pre expr}       & Precondici\'on \\
      \texttt{@post expr}      & Postcondici\'on \\
      \texttt{@invariant expr} & Invariante de bucle \\
      \texttt{@variant expr}   & Terminaci\'on \\
      \texttt{==>}             & Implicaci\'on \\
      \texttt{result}          & Valor de retorno \\
    \end{tabular}
  \end{alertblock}
  \begin{exampleblock}{Ejemplo}
    \small\texttt{@pre\; true}\\
    \texttt{@post result >= 0}\\
    \texttt{public static int abs()\{...\}}
  \end{exampleblock}
\end{columns}
\end{frame}

% ═════════════════════════════════════════════════════════════════════════════
\section{Front-end: L\'exico y Sint\'actico}
% ─────────────────────────────────────────────────────────────────────────────

\begin{frame}[fragile]{Analizador l\'exico (Flex)}
\begin{columns}[t]
  \column{0.46\textwidth}
  \begin{block}{Tokens principales}\small
    \begin{description}
      \item[\textbf{Reservadas}] exactas (\texttt{if, while, return,\ldots})
      \item[\textbf{Identificadores}] \texttt{[a-zA-Z][a-zA-Z0-9]*}
      \item[\textbf{Enteros}] \texttt{[0-9]+}
      \item[\textbf{Reales}] \texttt{[0-9]+\textbackslash.[0-9]+}
      \item[\textbf{Anotaciones}] \texttt{@pre}, \texttt{@post}, \texttt{@invariant}
    \end{description}
  \end{block}

  \column{0.54\textwidth}
  \begin{exampleblock}{Reglas especiales — Fase 5}
\begin{lstlisting}[style=cppslide]
"result"     { return RESULT_TK; }
"@pre"       { return PRE_TK; }
"@post"      { return POST_TK; }
"@invariant" { return INV_TK; }
"@variant"   { return VAR_TK; }
"==>"        { return IMPLIES_TK; }
/* debe ir antes de la regla de ID */
[a-zA-Z]... { return ID_TK; }
\end{lstlisting}
  \end{exampleblock}
\end{columns}
\end{frame}

% ─────────────────────────────────────────────────────────────────────────────
\begin{frame}[fragile]{Analizador sint\'actico (Bison LALR(1))}
\begin{columns}[t]
  \column{0.55\textwidth}
\begin{lstlisting}[style=cppslide]
Program : 'public' 'class' ID
          '{' Members '}'

Method  : 'public' 'static' Tipo
          ID '(' ')' Block

Instr : Asig | Return | If
      | 'while' '(' Expr ')'
        OptInv OptVar Block

/* Anotaciones formales */
OptPre  : '@pre'  Expr | /*empty*/
OptPost : '@post' Expr | /*empty*/
OptInv  : '@invariant' Expr | /*empty*/
OptVar  : '@variant'   Expr | /*empty*/
\end{lstlisting}

  \column{0.45\textwidth}
  \begin{block}{AST generado}
    \centering\small
    \begin{tikzpicture}[
      n/.style={rectangle,draw,rounded corners,
                font=\tiny\ttfamily,
                minimum width=2cm,
                minimum height=0.38cm,
                fill=blue!8, align=center},
      scale=0.9, transform shape
    ]
    \node[n] (prog) {ProgramNode};
    \node[n, below left=0.45cm and 0.6cm of prog]  (m1) {MethodNode};
    \node[n, below right=0.45cm and 0.4cm of prog] (m2) {MethodNode};
    \node[n, below=0.45cm of m1] (blk) {BlockNode};
    \node[n, below left=0.4cm and 0.1cm of blk]  (wh) {WhileNode};
    \node[n, below right=0.4cm and 0.1cm of blk] (ret) {ReturnNode};
    \draw[-Stealth] (prog)--(m1);
    \draw[-Stealth] (prog)--(m2);
    \draw[-Stealth] (m1)--(blk);
    \draw[-Stealth] (blk)--(wh);
    \draw[-Stealth] (blk)--(ret);
    \end{tikzpicture}
  \end{block}
  \begin{exampleblock}{\small Patr\'on Visitor}
    \texttt{SemanticVisitor}, \texttt{IRBuilder},\\
    \texttt{ASTVisualizer} — todos implementan\\
    la interfaz \texttt{Visitor} sin tocar el AST
  \end{exampleblock}
\end{columns}
\end{frame}

% ═════════════════════════════════════════════════════════════════════════════
\section{An\'alisis sem\'antico}
% ─────────────────────────────────────────────────────────────────────────────

\begin{frame}[fragile]{Tabla de s\'imbolos y comprobaciones}
\begin{columns}[t]
  \column{0.5\textwidth}
  \begin{block}{\texttt{TablaSimbolos}}
\begin{lstlisting}[style=cppslide]
struct Simbolo {
  string nombre;
  unsigned tipo;  // INT,REAL,BOOL
  unsigned dir;   // stack offset
  bool isConstant;
  bool isFunction;
  int  returnType;
};
class TablaSimbolos {
  TablaSimbolos* padre; // scope
  bool set(Simbolo s);
  Simbolo* get(string n);
};
\end{lstlisting}
  \end{block}

  \column{0.5\textwidth}
  \begin{block}{Comprobaciones}
    \begin{enumerate}\small
      \item Declaraci\'on antes de uso
      \item Compatibilidad de tipos
      \item No reasignaci\'on de \texttt{final}
      \item Tipo de retorno correcto
      \item C\'odigo inalcanzable
      \item Llamadas a funciones v\'alidas
    \end{enumerate}
  \end{block}
  \begin{exampleblock}{Coerci\'on de tipos}
    \small
    \texttt{int} op \texttt{double} $\Rightarrow$ \texttt{double}\\
    \texttt{int} / \texttt{int} $\Rightarrow$ \texttt{int}
  \end{exampleblock}
\end{columns}
\end{frame}

% ═════════════════════════════════════════════════════════════════════════════
\section{Representaci\'on Intermedia}
% ─────────────────────────────────────────────────────────────────────────────

\begin{frame}[fragile]{Cuadruplas de tres direcciones}
\begin{columns}[t]
  \column{0.46\textwidth}
  \begin{block}{Formato: $(op, arg_1, arg_2, res)$}\small
    \begin{description}\itemsep0pt
      \item[\texttt{IR\_ADD/SUB/MUL/DIV}] aritm\'etica
      \item[\texttt{IR\_MOD / IR\_NEG}] m\'odulo, negaci\'on
      \item[\texttt{IR\_COPY}] $x \leftarrow a$
      \item[\texttt{IR\_JZ / IR\_JMP}] saltos condicional/incond.
      \item[\texttt{IR\_LABEL}] etiqueta de destino
      \item[\texttt{IR\_CALL}] $t \leftarrow f()$
      \item[\texttt{IR\_RETURN}] devuelve valor
    \end{description}
  \end{block}

  \column{0.54\textwidth}
  \begin{exampleblock}{Ejemplo: while + suma}
\begin{lstlisting}[style=asmslide]
[FUNC_BEGIN] sumToFive
[LABEL]  L0
[LTE]    i 5    -> t0
[JZ]     t0     -> L1
[ADD]    s i    -> t1
[COPY]   t1     -> s
[ADD]    i 1    -> t2
[COPY]   t2     -> i
[JMP]    L0
[LABEL]  L1
[RETURN] s
[FUNC_END] sumToFive
\end{lstlisting}
  \end{exampleblock}
  \begin{block}{Optimizaciones ($-31\%$ instrucc.)}
    \small
    \textbf{1.} Elim. c\'odigo muerto \quad
    \textbf{2.} Plegado constantes \quad
    \textbf{3.} Prop. de copias
  \end{block}
\end{columns}
\end{frame}

% ═════════════════════════════════════════════════════════════════════════════
\section{Backend x86-64}
% ─────────────────────────────────────────────────────────────────────────────

\begin{frame}[fragile]{Generaci\'on de c\'odigo x86-64}
\begin{columns}[t]
  \column{0.5\textwidth}
  \begin{block}{ABI System V AMD64}\small
    \begin{description}\itemsep0pt
      \item[Argumentos] \reg{rdi}, \reg{rsi}, \reg{rdx},\ldots\reg{r9}
      \item[Retorno int] \reg{rax}
      \item[Retorno float] \reg{xmm0}
      \item[Caller-saved] \reg{rcx}, \reg{rdx}, \reg{r8}\ldots\reg{r11}
    \end{description}
  \end{block}
\begin{lstlisting}[style=asmslide]
; prologo / epilogo estandar
pushq  %rbp
movq   %rsp, %rbp
subq   $N, %rsp   ; N % 16 == 0
...
movq   %rbp, %rsp
popq   %rbp
ret
\end{lstlisting}

  \column{0.5\textwidth}
  \begin{exampleblock}{Bucle compilado}
\begin{lstlisting}[style=asmslide]
sumToFive:
    pushq  %rbp
    movq   %rsp, %rbp
    subq   $48, %rsp
L0:
    movq   -8(%rbp), %rax  ; i
    cmpq   $5, %rax
    setle  %al
    movzbq %al, %rax
    movq   %rax, -24(%rbp) ; t0
    testq  %rax, %rax
    je     L1
    ; s = s + i
    movq  -16(%rbp), %rax
    addq  -8(%rbp),  %rax
    movq  %rax, -16(%rbp)
    jmp   L0
L1: movq  -16(%rbp), %rax
    ...
    ret
\end{lstlisting}
  \end{exampleblock}
\end{columns}
\end{frame}

% ═════════════════════════════════════════════════════════════════════════════
\section{Asignaci\'on de registros}
% ─────────────────────────────────────────────────────────────────────────────

\begin{frame}[fragile]{Chaitin-Briggs: coloraci\'on de grafos ($K=6$)}
\begin{columns}[t]
  \column{0.55\textwidth}
  \begin{block}{Tres fases}
    \begin{enumerate}\small
      \item \textbf{Vivacidad} (punto fijo)
        \[ \text{live\_in}[B] = \text{use}[B] \cup (\text{live\_out}[B] \setminus \text{def}[B]) \]
      \item \textbf{Grafo de interferencias}\\
        Arista $(d,v)$ si $d$ se define con $v$ viva
      \item \textbf{Coloraci\'on}
        \begin{itemize}\small
          \item \emph{Simplify}: apila nodos grado $<K$
          \item \emph{Spill}: derrama el de mayor grado
          \item \emph{Assign}: asigna colores al desapilar
        \end{itemize}
    \end{enumerate}
  \end{block}

  \column{0.45\textwidth}
  \begin{block}{Registros disponibles}
    \small
    \begin{tabular}{ll}
      Enteros: & \reg{rcx}, \reg{rdx}, \reg{r8}\\
               & \reg{r9}, \reg{r10}, \reg{r11}\\
      Floats:  & \reg{xmm2}\ldots\reg{xmm7}\\
    \end{tabular}
  \end{block}
  \begin{exampleblock}{Resultado}
\begin{lstlisting}[style=asmslide]
$ ./compiler --reg-alloc t04.java
  i  -> %rcx
  s  -> %rdx
  t0 -> %r8
  t1 -> spill (-32(%rbp))
\end{lstlisting}
  \end{exampleblock}
\end{columns}
\end{frame}

% ═════════════════════════════════════════════════════════════════════════════
\section{Verificaci\'on formal}
% ─────────────────────────────────────────────────────────────────────────────

\begin{frame}{L\'ogica de Hoare y c\'alculo WP}
\begin{columns}[t]
  \column{0.48\textwidth}
  \begin{block}{Tripleta de Hoare}
    \[ \hoare{P}{S}{Q} \]
    Si $P$ vale antes de $S$ y termina, $Q$ vale despu\'es.
    \medskip

    \textbf{Condici\'on principal:}
    \[ P \Rightarrow \WP{\mathit{body}}{Q} \]
    v\'alida $\iff$ Z3 dice \texttt{UNSAT} a $\neg\mathit{VC}$
  \end{block}

  \column{0.52\textwidth}
  \begin{block}{Reglas del c\'alculo WP}
    \begin{align*}
      \WP{x:=e}{Q} &= \subst{Q}{x}{e}\\[3pt]
      \WP{S_1;S_2}{Q} &= \WP{S_1}{\WP{S_2}{Q}}\\[3pt]
      \WP{\texttt{ret}\;e}{Q} &= \subst{Q}{\texttt{result}}{e}\\[3pt]
      \WP{\texttt{if}(B)\;S_1\;\texttt{else}\;S_2}{Q}
        &= (B \Rightarrow \WP{S_1}{Q})\\
        &\;\wedge (\neg B \Rightarrow \WP{S_2}{Q})
    \end{align*}
  \end{block}
\end{columns}
\end{frame}

% ─────────────────────────────────────────────────────────────────────────────
\begin{frame}[fragile]{Implementaci\'on CPS para retornos anticipados}
\begin{alertblock}{Problema con la regla na\"ive}
Para \;\texttt{if(x<0) return -x; \; return x;}\; con $Q=\texttt{result}\geq 0$:
la composici\'on inversa pierde \texttt{result} tras el primer \texttt{return}.
\end{alertblock}

\begin{exampleblock}{Soluci\'on: continuaciones (procesado hacia adelante)}
\begin{lstlisting}[style=cppslide]
Predicate* wpSeq(stmts, idx, Q) {
    auto* s = stmts[idx];
    // return: usa Q directamente -- ignora continuacion
    if (auto* r = dyn_cast<ReturnNode>(s))
        return wpReturn(r, Q);
    // if: la continuacion se propaga a ambas ramas
    if (auto* i = dyn_cast<IfNode>(s)) {
        vector<StmtNode*> cont(stmts.begin()+idx+1, stmts.end());
        return And(B ==> wpBranchAndCont(then, cont, Q),
                  !B ==> wpBranchAndCont(else, cont, Q));
    }
    // assign: WP(x:=e; resto, Q) = WP(resto, Q)[e/x]
    if (auto* a = dyn_cast<AssignNode>(s))
        return wpAssign(a, wpSeq(stmts, idx+1, Q));
    return wpSeq(stmts, idx+1, Q);
}
\end{lstlisting}
\end{exampleblock}
\end{frame}

% ─────────────────────────────────────────────────────────────────────────────
\begin{frame}{VCs de bucles con invariante}
Para \texttt{while($B$)\{$S$\}} con invariante $I$:

\begin{columns}[t]
  \column{0.5\textwidth}
  \begin{block}{Tres condiciones de verificaci\'on}
    \medskip
    $\mathit{VC}_1$ \textbf{(iniciaci\'on)}:
    \[ P \Rightarrow \WP{S_{\text{antes}}}{I} \]
    $\mathit{VC}_2$ \textbf{(preservaci\'on)}:
    \[ (I \wedge B) \Rightarrow \WP{S_{\text{body}}}{I} \]
    $\mathit{VC}_3$ \textbf{(uso)}:
    \[ (I \wedge \neg B) \Rightarrow \WP{S_{\text{desp}}}{Q} \]
  \end{block}

  \column{0.5\textwidth}
  \begin{alertblock}{El invariante debe ser \emph{fuerte}}
    Para la suma $s = 1+\cdots+i$:
    \medskip

    $I = s\geq 0$ \quad insuficiente:\\
    \quad Z3 halla $(i=-1,\,s=0)$ que satisface $I\wedge B$\\
    \quad pero $(s+i)=s-1<0$. \quad \textcolor{warn}{\textbf{Fallo.}}
    \medskip

    $I = s\geq 0 \;\wedge\; i\geq 1$ \quad correcto:
    \[ (s\geq0 \wedge i\geq1 \wedge i\leq5) \Rightarrow (s+i)\geq0 \]
    Z3: \texttt{UNSAT} \quad\textcolor{accent}{$\checkmark$}
  \end{alertblock}
\end{columns}
\end{frame}

% ─────────────────────────────────────────────────────────────────────────────
\begin{frame}[fragile]{SMT-Lib2 y Z3}
\begin{columns}[t]
  \column{0.55\textwidth}
\begin{lstlisting}[style=smtslide]
(push 1)
; [absValue] correccion
(declare-fun x () Int)
(assert (not
  (=> true
    (and
      (=> (< (- 5) 0)
          (>= (- (- 5)) 0))
      (=> (not (< (- 5) 0))
          (>= (- 5) 0))
    )
  )
))
(check-sat)   ; --> unsat  [OK]
(pop 1)
\end{lstlisting}

  \column{0.45\textwidth}
  \begin{block}{Interpretaci\'on de Z3}\small
    \begin{description}\itemsep2pt
      \item[\texttt{unsat}] VC v\'alida \quad {\color{accent}\textbf{[OK]}}
      \item[\texttt{sat}] Bug hallado \quad {\color{warn}\textbf{[BUG]}}
      \item[\texttt{unknown}] Timeout \quad \textbf{[?]}
    \end{description}
  \end{block}
  \begin{block}{Tipos de VC generadas}\small
    \begin{enumerate}
      \item Correcci\'on del m\'etodo
      \item Iniciaci\'on del invariante
      \item Preservaci\'on del invariante
      \item Uso del invariante
      \item Terminaci\'on (\texttt{@variant})
      \item Seguridad: divisi\'on por cero
    \end{enumerate}
  \end{block}
\end{columns}
\end{frame}

% ─────────────────────────────────────────────────────────────────────────────
\begin{frame}[fragile]{Demo: verificaci\'on completa}
\begin{columns}[t]
  \column{0.52\textwidth}
\begin{lstlisting}[style=javaslide]
@pre  true
@post result >= 0
public static int absValue() {
    int x; x = -5;
    if (x < 0) return -x;
    return x;
}

@post result > 10          // BUG!
public static int wrongMethod() {
    int x; x = 5; return x;
}

@post result >= 0
public static int sumLoop() {
    int i; int s; i=1; s=0;
    while (i <= 5)
    @invariant s>=0 && i>=1
    @variant   6 - i
    { s=s+i; i=i+1; }
    return s;
}
\end{lstlisting}

  \column{0.48\textwidth}
  \begin{exampleblock}{Salida del verificador}
\begin{lstlisting}[style=asmslide]
$ ./compiler --verify demo.java

[OK]  absValue
      correccion (unsat)

[OK]  sumLoop
      iniciacion (unsat)
[OK]  sumLoop
      preservacion (unsat)
[OK]  sumLoop
      uso (unsat)
[OK]  sumLoop
      correccion (unsat)
[OK]  sumLoop
      terminacion (unsat)

[BUG] wrongMethod
      correccion (sat) <-- bug!
\end{lstlisting}
  \end{exampleblock}
  \bigskip
  \centering\textbf{6 OK} $\;\cdot\;$ \textcolor{warn}{\textbf{1 BUG detectado}}
\end{columns}
\end{frame}

% ═════════════════════════════════════════════════════════════════════════════
\section{Pruebas y resultados}
% ─────────────────────────────────────────────────────────────────────────────

\begin{frame}{Suite de pruebas: 11/11}
\begin{columns}[t]
  \column{0.33\textwidth}
  \begin{block}{Backend (5)}
    \begin{tabular}{ll}\small
      t01\_print    & {\color{accent}\texttt{PASS}} \\
      t02\_arith    & {\color{accent}\texttt{PASS}} \\
      t03\_if\_else & {\color{accent}\texttt{PASS}} \\
      t04\_while    & {\color{accent}\texttt{PASS}} \\
      t05\_funcs    & {\color{accent}\texttt{PASS}} \\
    \end{tabular}
  \end{block}

  \column{0.33\textwidth}
  \begin{block}{Sem\'antico (3)}
    \begin{tabular}{ll}\small
      sem01\_undecl & {\color{accent}\texttt{PASS}} \\
      sem02\_type   & {\color{accent}\texttt{PASS}} \\
      sem03\_redecl & {\color{accent}\texttt{PASS}} \\
    \end{tabular}
  \end{block}

  \column{0.34\textwidth}
  \begin{block}{Verificaci\'on (3)}
    \begin{tabular}{ll}\small
      ver01\_simple & {\color{accent}\texttt{PASS}} \\
      ver02\_inv    & {\color{accent}\texttt{PASS}} \\
      demo\_verified & {\color{accent}\texttt{PASS}} \\
    \end{tabular}
  \end{block}
\end{columns}

\vspace{0.4cm}
\begin{beamercolorbox}[rounded=true, center, wd=0.6\textwidth]{block title}
  \centering\large\bfseries
  \texttt{make \&\& tests/run\_tests.sh} $\Rightarrow$ \textbf{11 / 11 passed}
\end{beamercolorbox}
\end{frame}

% ═════════════════════════════════════════════════════════════════════════════
\section{Conclusiones}
% ─────────────────────────────────────────────────────────────────────────────

\begin{frame}{Resumen de milestones}
\begin{table}
  \begin{tabular}{clll}
    \hline
    \textbf{Fase} & \textbf{Milestone} & \textbf{Componentes clave} & \\
    \hline
    1 & Front-end  & Flex + Bison, AST C++ & {\color{accent}$\checkmark$} \\
    2 & Sem\'antico & \texttt{TablaSimbolos}, tipos, Visitor & {\color{accent}$\checkmark$} \\
    3 & IR + Optim. & Cuadruplas, dead code, folding, copy & {\color{accent}$\checkmark$} \\
    4 & Backend    & x86-64, ABI, Chaitin-Briggs & {\color{accent}$\checkmark$} \\
    5 & Verificaci\'on & Hoare, WP-CPS, Z3, VCGenerator & {\color{accent}$\checkmark$} \\
    \hline
  \end{tabular}
\end{table}

\begin{columns}[t]
  \column{0.5\textwidth}
  \begin{block}{Aportaciones principales}
    \begin{itemize}\small
      \item Compilador completo end-to-end
      \item WP con continuaciones (CPS) para retornos anticipados
      \item VCs autom\'aticas: correcci\'on, invariantes, terminaci\'on
      \item Integraci\'on Z3 con reporte coloreado
    \end{itemize}
  \end{block}

  \column{0.5\textwidth}
  \begin{alertblock}{Trabajo futuro}
    \begin{itemize}\small
      \item Inferencia autom\'atica de invariantes (Daikon)
      \item Soporte de arrays y objetos
      \item SSA form para optimizaci\'on avanzada
      \item Backend LLVM-IR
    \end{itemize}
  \end{alertblock}
\end{columns}
\end{frame}

% ─── Diapositiva final ────────────────────────────────────────────────────────
\begin{frame}[plain]
  \centering
  \vspace{1.5cm}
  {\Huge\bfseries\color{primary} Gracias}\\[0.8em]
  {\large Compilador \texttt{JavaSubset} $\to$ x86-64}\\[0.4em]
  {\normalsize con Verificaci\'on Formal mediante Hoare + Z3}\\[1.5em]
  \hrule
  \vspace{0.6em}
  \begin{tabular}{rl}
    \texttt{make \&\& tests/run\_tests.sh} & $\Rightarrow$ \textbf{11/11 passed}\\
    \texttt{./compiler --verify demo.java} & $\Rightarrow$ \textbf{6 OK, 1 BUG}\\
    \texttt{pdflatex presentation.tex}     & $\Rightarrow$ \textbf{estas diapositivas}\\
  \end{tabular}
\end{frame}

\end{document}
